\chapter{\ifenglish Conclusions and Discussions\else บทสรุปและข้อเสนอแนะ\fi}

ในบทนี้จะนำเสนอการสรุปภาพรวมของโครงงานที่ได้ดำเนินงานมาทั้งหมด รวมถึงปัญหาที่พบในระหว่างกระบวนการวิจัยและพัฒนา ตลอดจนข้อเสนอแนะสำหรับการต่อยอดในอนาคต เพื่อประโยชน์แก่ผู้ที่สนใจศึกษาในด้านความมั่นคงปลอดภัยไซเบอร์เชิงควอนตัมต่อไป

\section{สรุปผลการดำเนินงาน (Project Summary)}
โครงงานนี้ได้สำเร็จลุล่วงในการออกแบบและสร้างแบบจำลองทางคณิตศาสตร์ในรูปแบบ QUBO เพื่อค้นหาดุลยภาพแนช (Nash Equilibrium) ในสถานการณ์จำลองความมั่นคงปลอดภัยทางไซเบอร์ โดยสามารถสรุปประเด็นสำคัญได้ดังนี้:

\begin{itemize}
    \item \textbf{การพัฒนาแบบจำลอง:} ผู้วิจัยสามารถแปลงโจทย์ความปลอดภัยไซเบอร์ให้อยู่ในรูปโครงสร้างต้นไม้ (Cyber Tree) และประเมินความเสี่ยงตามมาตรฐาน ISO/IEC 27001 จนสามารถสร้างตารางผลตอบแทน (Payoff Matrix) ที่ครอบคลุมทั้งฝ่ายโจมตีและฝ่ายป้องกันได้สำเร็จ
    \item \textbf{การเข้ารหัส QUBO:} ได้ออกแบบสมการ QUBO ที่มีประสิทธิภาพโดยใช้เทคนิค Slack Variable และ Binary Expansion เพื่อจำลองเงื่อนไขความเสียใจ (Regret) ทำให้คอมพิวเตอร์ควอนตัมสามารถแยกแยะสถานะที่เป็นดุลยภาพแนชออกจากสถานะอื่นได้อย่างชัดเจน
    \item \textbf{อัลกอริทึมแบบไฮบริด (Hybrid Algorithm):} ประสบความสำเร็จในการนำแนวคิด Hybrid Quantum-Classical มาประยุกต์ใช้เพื่อแก้ไขข้อจำกัดด้านทรัพยากรคิวบิต (Qubit limitations) โดยให้คอมพิวเตอร์แบบดั้งเดิมจัดการการวนรอบกลยุทธ์ และส่งเฉพาะส่วนคำนวณเชิงลึกให้เครื่อง Quantum Annealer ทำให้สามารถแก้ปัญหาที่มีขนาดใหญ่กว่าการประมวลผลบนควอนตัมเพียงอย่างเดียวได้
    \item \textbf{ความถูกต้องของผลลัพธ์:} จากการทดสอบทั้งในกรณีผู้เล่น 2 คน และ 3 คน (2 Attackers vs 1 Defender) ผลการรันแสดงให้เห็นว่าระบบสามารถระบุจุดสมดุลที่ค่าพลังงานต่ำสุด (Energy = 0) ได้อย่างแม่นยำ สอดคล้องกับหลักฐานการพิสูจน์ทางคณิตศาสตร์
\end{itemize}

\section{ปัญหาที่พบและแนวทางการแก้ไข (Problems and Solutions)}
ในการดำเนินงานมีข้อจำกัดและปัญหาเชิงเทคนิคที่สำคัญ ซึ่งได้ดำเนินการแก้ไขและหาแนวทางรองรับไว้ดังนี้:

\begin{enumerate}
    \item \textbf{ข้อจำกัดด้านทรัพยากรคิวบิต:} หากประมวลผลตารางผลตอบแทนขนาดใหญ่พร้อมกันทั้งหมด จะใช้จำนวนคิวบิตเกินกว่าที่ระบบรองรับ
    \begin{itemize}
        \item \textbf{การแก้ไข:} ใช้เทคนิค Qubit Reuse ผ่านระบบ Hybrid โดยการตรวจสอบทีละเซลล์กลยุทธ์ (Cell-based verification) แทนการแก้ปัญหาทั้งระบบในครั้งเดียว
    \end{itemize}
    \item \textbf{ความแม่นยำในการจำลอง (Simulation Accuracy):} ในกรณีปัญหาที่มีความซับซ้อนสูงหรือมีขนาดใหญ่ขึ้น คำตอบจากเครื่องควอนตัมอาจมีความคลาดเคลื่อนเล็กน้อยจากสภาวะพลังงานต่ำสุดเฉพาะจุด (Local Minima)
    \item \textbf{การกำหนด Penalty Coefficient:} หากค่าพารามิเตอร์บทลงโทษไม่เหมาะสม อาจทำให้ระบบหาจุดสมดุลผิดพลาด
    \begin{itemize}
        \item \textbf{การแก้ไข:} มีการทดลองปรับค่าพารามิเตอร์ $P$ และ $W$ จนได้สัดส่วนที่เหมาะสมที่ทำให้ระบบหลีกเลี่ยงสถานะที่ไม่ใช่แนชได้อย่างสมบูรณ์
    \end{itemize}
\end{enumerate}

\section{ข้อเสนอแนะและแนวทางการพัฒนาต่อ (Future Work and Recommendations)}
เพื่อให้แบบจำลองนี้มีความสมบูรณ์และสามารถนำไปใช้งานได้จริงในวงกว้าง ผู้วิจัยมีข้อเสนอแนะดังนี้:

\begin{itemize}
    \item \textbf{การวิเคราะห์ความลึกของดุลยภาพ (The Best of Nash):} พัฒนาอัลกอริทึมเพิ่มเติมในการเปรียบเทียบจุดสมดุลแนชในกรณีที่ระบบตรวจพบหลายจุด (Multiple Equilibria) โดยใช้การวิเคราะห์เพื่อนบ้าน (Neighbor analysis) เพื่อค้นหาจุดสมดุลที่มีประสิทธิภาพรวมสูงที่สุด
    \item \textbf{การเพิ่มความยืดหยุ่นของต้นไม้ (Dynamic Cyber Tree):} พัฒนาให้แบบจำลองสามารถรองรับโครงสร้างต้นไม้ที่มีความเปลี่ยนแปลงได้ตลอดเวลา (Dynamic Tree) เพื่อจำลองเครือข่ายคอมพิวเตอร์ที่ใช้งานจริงซึ่งมีการเชื่อมต่ออุปกรณ์ใหม่ ๆ อยู่เสมอ
    \item \textbf{การรวมเทคโนโลยี AI/ML:} นำแนวคิด Heuristic หรือ Machine Learning มาใช้ร่วมกับ Quantum Computing เพื่อช่วยในการคาดการณ์พฤติกรรมผู้โจมตีที่ซับซ้อนขึ้น และช่วยในการเลือกตัวแปรนำเข้าสำหรับ QUBO โดยอัตโนมัติ
    \item \textbf{การประมวลผลบนเครื่องควอนตัมจริง (Real QPU):} ขยายขอบเขตการรันงานจากเครื่องจำลอง (Simulator) ไปสู่การรันบนตัวประมวลผลควอนตัมจริงในระดับที่ใหญ่ขึ้น เพื่อทดสอบประสิทธิภาพในเชิงเวลา (Time Complexity) ต่อไป
\end{itemize}